problem:  
Let \(M\) be a fixed smooth, compact \(2n\)-dimensional manifold that admits at least one complex structure \(J\) for which \((M,J)\) is a smooth Fano manifold.  
For any such complex structure \(J\), denote by \(\omega_J\) the unique Kähler–Einstein metric satisfying \(\mathrm{Ric}(\omega_J)=\omega_J\), viewed as a Riemannian metric on the underlying smooth manifold \(M\).

**Question:**  
Can there exist two *non-biholomorphic* complex structures \(J_1, J_2\) on \(M\) such that:
1. each \((M,J_i)\) is a smooth Fano manifold admitting a Kähler–Einstein metric \(\omega_{J_i}\), and  
2. the Riemannian metric spaces \((M,\omega_{J_1})\) and \((M,\omega_{J_2})\) are *isometric* as real Riemannian manifolds?  

Equivalently, does the Riemannian isometry class of the Kähler–Einstein metric on a Fano manifold uniquely determine its complex structure up to biholomorphism?

If such examples exist, must they arise from Fano manifolds with continuous automorphism groups, or can this phenomenon occur even when each \(\mathrm{Aut}(M,J_i)\) is discrete?  
Conversely, if no such examples exist, does this imply a *metric Torelli property* for Kähler–Einstein Fano manifolds—that the real Riemannian metric suffices to reconstruct the complex structure uniquely (up to holomorphic isomorphism)?

---

Why is it a "good" problem:  
This question distills a subtle rigidity issue at the intersection of Riemannian and complex geometry: whether the analytic “shape” of a Kähler–Einstein Fano manifold determines its complex structure uniquely. It is precise, simple to state, yet remains open even in basic cases—no known classification or obstruction settles the possibility of distinct complex Fano structures sharing an identical Einstein metric.

The problem unites several central themes:
- **Metric rigidity and automorphism theory:** An isometry between two Kähler–Einstein metrics need not preserve complex structures, raising the question of when holomorphic and isometric symmetries differ.
- **Moduli interactions:** If injectivity holds, the metric moduli (in Gromov–Hausdorff or Cheeger–Gromov sense) would faithfully reflect complex moduli, strengthening the analytic–algebraic link predicted by the Yau–Tian–Donaldson correspondence.
- **Potential counterexamples:** Existence of such pairs would reveal hidden degeneracies in the relationship between KE metrics and complex structure, showing that identical Riemannian geometry can encode distinct complex geometries.

In essence, the problem asks whether the Kähler–Einstein equation “remembers” the complex structure entirely, or whether distinct algebraic geometries can manifest the same Einstein Riemannian universe—a question that is conceptually minimal yet deeply foundational across complex differential geometry, algebraic moduli theory, and metric analysis.